\section{Knowledge Extraction}
\label{sec:knowledge}

\subsection{Fine-Grained Knowledge Extraction (Task 6)}

For each video section, we extract structured knowledge points using a large language model. This transforms unstructured transcript text into organized, searchable knowledge entries.

\subsubsection{LLM Configuration}

We use Qwen2.5-7b-instruct via Alibaba DashScope's OpenAI-compatible API:

\begin{itemize}
    \item \textbf{Model:} qwen2.5-7b-instruct (7 billion parameters)
    \item \textbf{Temperature:} 0.3 (low for consistent structured output)
    \item \textbf{Max tokens:} 2048
    \item \textbf{Response format:} JSON
\end{itemize}

\subsubsection{Prompt Template}

The extraction prompt is carefully engineered to produce consistent, high-quality output:

\begin{lstlisting}[language=text, caption={Fine-Grained Extraction Prompt}]
You are an expert educational content analyst. Analyze the 
following lecture segment and extract structured knowledge points.

Section: {section_title}
Time range: {time_range}

Transcript:
---
{transcript}
---

Extract the following in JSON format:
{
  "section_title": "A concise, descriptive title (5-10 words)",
  "points": [
    {
      "title": "Knowledge point title (3-8 words)",
      "summary": "2-3 sentence explanation",
      "terms": ["key term 1", "key term 2"],
      "importance": 0.0-1.0
    }
  ]
}

Rules:
- Extract 1-5 knowledge points per segment
- Each point should be self-contained and understandable
- Key terms should be specific technical vocabulary
- Importance reflects how central the concept is to the lecture
\end{lstlisting}

\subsubsection{Response Parsing}

LLM responses are parsed with robust error handling:

\begin{enumerate}
    \item Remove markdown code fences (\texttt{\`\`\`json ... \`\`\`})
    \item Extract JSON from prose if embedded in text
    \item Validate against expected schema
    \item Handle parsing errors gracefully (log and continue)
\end{enumerate}

\subsubsection{Knowledge Point Schema}

Each extracted knowledge point follows this structure:

\begin{lstlisting}[language=json, caption={Knowledge Point Structure}]
{
  "section_title": "Gradient Descent Optimization",
  "points": [
    {
      "title": "Learning Rate Parameter",
      "summary": "The learning rate (alpha) controls the step size 
                  during gradient descent. Too large causes divergence, 
                  too small results in slow convergence.",
      "terms": ["learning rate", "gradient descent", "convergence"],
      "importance": 0.85
    },
    {
      "title": "Cost Function Minimization",
      "summary": "Gradient descent iteratively minimizes the cost 
                  function J(theta) by moving in the direction of 
                  steepest descent (negative gradient).",
      "terms": ["cost function", "gradient", "minimization"],
      "importance": 0.92
    }
  ]
}
\end{lstlisting}

\subsubsection{Error Isolation}

To ensure robustness, the task implements error isolation at the section level:

\begin{itemize}
    \item Sections with transcript text $< 20$ characters are skipped (likely silence-only segments)
    \item LLM call failures are logged but don't halt processing
    \item JSON parsing errors are caught and logged for debugging
    \item The task continues to subsequent sections even if some fail
\end{itemize}

This design ensures that issues in one section don't prevent knowledge extraction from the entire video.

\subsection{Coarse-Grained Knowledge Summarization (Task 8)}

After fine-grained extraction, all knowledge points are aggregated into a lecture-level summary.

\subsubsection{Aggregation Process}

The coarse-grained task gathers all fine-grained knowledge:

\begin{lstlisting}[language=json, caption={Input to Coarse-Grained Task}]
{
  "sections": [
    {
      "title": "Introduction to Gradient Descent",
      "time_range": "0:00-5:30",
      "points": [
        {"title": "...", "summary": "...", "terms": [...]},
        ...
      ]
    },
    ...
  ]
}
\end{lstlisting}

This structured input is sent to the LLM with an aggregation prompt requesting:

\begin{itemize}
    \item Overall lecture summary (3-5 paragraphs)
    \item Thematic chapter outline with time ranges
    \item Key concepts list (ordered from foundational to advanced)
    \item Learning objectives (actionable and measurable)
    \item Difficulty assessment (beginner/intermediate/advanced)
\end{itemize}

\subsubsection{Output Structure}

\begin{lstlisting}[language=json, caption={Coarse-Grained Summary Output}]
{
  "summary": "This lecture introduces gradient descent, a fundamental 
              optimization algorithm in machine learning...",
  "chapters": [
    {
      "title": "Introduction and Motivation",
      "begin_time": 0.0,
      "end_time": 300.0,
      "summary": "Overview of optimization in ML",
      "key_points": ["What is optimization", "Why gradient descent"]
    },
    ...
  ],
  "concepts": ["Gradient Descent", "Learning Rate", "Cost Function", 
               "Convergence", "Batch vs Stochastic"],
  "objectives": [
    "Explain the gradient descent algorithm",
    " tune the learning rate parameter",
    "Compare batch and stochastic gradient descent"
  ],
  "difficulty": "intermediate"
}
\end{lstlisting}

\subsection{Mindmap Generation (Task 9)}

The final task generates a hierarchical concept map for visualization.

\subsubsection{Tree Structure}

The mindmap is represented as a nested tree structure:

\begin{lstlisting}[language=json, caption={Mindmap JSON Structure}]
{
  "root": {
    "id": "root",
    "label": "Machine Learning Fundamentals",
    "children": [
      {
        "id": "chapter1",
        "label": "Supervised Learning",
        "time_range": [0, 600],
        "children": [
          {
            "id": "kp1",
            "label": "Linear Regression",
            "time_range": [60, 180]
          },
          {
            "id": "kp2",
            "label": "Classification",
            "time_range": [180, 360]
          }
        ]
      }
    ]
  }
}
\end{lstlisting}

\subsubsection{Visualization}

The frontend renders this structure using React Flow, an interactive graph visualization library. Features include:

\begin{itemize}
    \item Collapsible nodes for large concept hierarchies
    \item Click-to-seek: clicking a node jumps to the relevant video timestamp
    \item Zoom and pan for navigating large mindmaps
    \item Color coding by chapter/theme
    \item Export to PNG/SVG for study materials
\end{itemize}
